\documentclass[12pt,a4paper]{article}

\usepackage[utf8]{inputenc}
\usepackage[T1]{fontenc}
\usepackage[spanish,es-noshorthands]{babel}
\usepackage[left=3cm,right=3cm,top=3cm,bottom=3cm]{geometry}
\usepackage{setspace}
\usepackage{graphicx}
\usepackage{booktabs}
\usepackage{multirow}
\usepackage{amsmath}
\usepackage{amssymb}
\usepackage{float}
\usepackage{caption}
\usepackage{subcaption}
\usepackage{enumitem}
\usepackage[round,authoryear]{natbib}
\usepackage{hyperref}

% Babel en espanol rotula las tablas como "Cuadro" y el indice como
% "Indice"; el formato del curso usa "Tabla" y "Tabla de contenidos".
% Babel reasigna estos nombres al abrir el documento, de modo que hay que
% engancharlos a captionsspanish y no redefinirlos sueltos en el preambulo.
\addto\captionsspanish{%
  \renewcommand{\tablename}{Tabla}%
  \renewcommand{\contentsname}{Tabla de contenidos}%
  \renewcommand{\refname}{Referencias}%
  \renewcommand{\bibname}{Referencias}%
}

\onehalfspacing
\setlength{\parindent}{1.2em}
\captionsetup{font=small,labelfont=bf}

\newcommand{\figdir}{../../reports/figures}

\begin{document}

% ============================================================ PORTADA
\begin{titlepage}
\thispagestyle{empty}
\begin{minipage}{0.75\textwidth}
  \raggedright
  {\large\bfseries UNIVERSIDAD DE SANTIAGO DE CHILE}\\[0.5cm]
  {\bfseries FACULTAD DE INGENIERÍA}\\[0.5cm]
  {\bfseries Departamento de Ingeniería Informática}
\end{minipage}
\hfill
\begin{minipage}{0.2\textwidth}
  \raggedleft
  \includegraphics[width=2.4cm]{img/escudo_usach.png}
\end{minipage}

\vspace{2.2cm}

\begin{center}
{\LARGE\bfseries Laboratorio 1: Pipeline de Aprendizaje Profundo\\[0.35cm]
sobre Datos Masivos\par}

\vspace{0.7cm}

{\Large\bfseries Clasificación de eventos de bosón de Higgs\\[0.35cm]
mediante Bagging y Stacking Ensemble\par}

\vspace{1.8cm}

{\large\bfseries Cristopher René Alexander Angulo Ahumada\par}
\end{center}

\vspace{1.8cm}

\begin{flushright}
\textbf{Profesor:} Francisco Muñoz\\[0.25cm]
\textbf{Curso:} Fundamentos de Aprendizaje Profundo y Aplicaciones\\[0.25cm]
\textbf{Programa:} Magíster en Informática\\[0.25cm]
\textbf{Grupo:} 6 --- Trabajo individual
\end{flushright}

\vfill

\begin{center}
\textbf{Santiago -- Chile}\\[0.25cm]
\textbf{2026}
\end{center}
\end{titlepage}

% ============================================================ INDICE
\newpage
\thispagestyle{empty}
\tableofcontents
\newpage
\setcounter{page}{1}

% ============================================================ CUERPO
\section{Descripción del Problema}

El \emph{HIGGS Boson Dataset} \citep{baldi2014, uci} contiene 11.000.000 de colisiones simuladas
de partículas en un acelerador. La tarea es de clasificación binaria: distinguir un evento de
\textbf{señal} ---producción de un bosón de Higgs--- de uno de \textbf{fondo}, entendido como un
proceso físico distinto que genera exactamente los mismos productos de decaimiento observables.

El problema reúne las tres condiciones que justifican el uso de aprendizaje profundo. Las
variables son continuas y están fuertemente correlacionadas entre sí; las diferencias entre ambas
clases son sutiles y no se manifiestan en ninguna variable aislada; y el volumen disponible
permite ajustar modelos de alta capacidad sin incurrir de inmediato en sobreajuste. Es, además, el
escenario en el que el trabajo original demostró que una red profunda supera a los métodos
tradicionales de la física de altas energías.

\subsection{Estructura de las variables}

Las 28 variables predictoras se organizan en dos niveles. Las 21 \textbf{de bajo nivel} son
mediciones directas del detector: para cada objeto reconstruido ---el leptón, la energía
transversal faltante y hasta cuatro \emph{jets}--- se registra su momento transversal $p_T$, su
pseudorapidez $\eta$, su ángulo azimutal $\phi$ y, en el caso de los \emph{jets}, su etiquetado
$b$. Las 7 \textbf{de alto nivel} son masas invariantes construidas por especialistas del dominio
a partir de las anteriores, con el fin explícito de discriminar señal de fondo; entre ellas
\texttt{m\_bb}, que en los eventos de señal reconstruye el bosón de Higgs a partir del par de
\emph{jets} $b$.

\subsection{Cumplimiento de los requisitos}

Se extrae una submuestra aleatoria de \textbf{1.000.000 de filas}, dentro del rango de 200.000 a
1.000.000 exigido. Respecto de la dimensionalidad, el conjunto aporta 28 columnas y su uso fue
autorizado expresamente por el profesor dado el volumen de observaciones disponible.
Adicionalmente, la ingeniería de variables de la Fase~B eleva el conjunto a \textbf{51
características}, con lo que el mínimo de 40 columnas queda satisfecho.

\section{Fase A: Carga Masiva, Particiones y Análisis Exploratorio}

\subsection{Gestión de memoria}

El archivo descomprimido ocupa 7,5\,GB. Una lectura directa con \texttt{pd.read\_csv()} intentaría
materializar las 11.000.000 de filas en memoria y, con la descompresión del \texttt{.gz}
ocurriendo en el mismo proceso, el desenlace habitual es la terminación de la sesión por
agotamiento de memoria. La estrategia adoptada combina cuatro medidas: lectura por bloques de
1.000.000 de filas mediante \texttt{chunksize}, casteo explícito a \texttt{float32} ---la mitad de
memoria por valor---, muestreo aleatorio del 9,09\,\% de \emph{cada} bloque y liberación explícita
de cada bloque mediante \texttt{del} y \texttt{gc.collect()}.

El muestreo uniforme sobre los 11 bloques no es un detalle accesorio: el orden de las filas en
HIGGS no es aleatorio, puesto que el \emph{benchmark} original reserva las últimas 500.000
observaciones como conjunto de prueba. Leer secuencialmente desde el inicio, o muestrear solo los
primeros bloques, habría producido particiones sesgadas y \emph{data drift} efectivo.

El consumo máximo del proceso fue de \textbf{0,94\,GB}, menos del 13\,\% del tamaño del archivo.
La progresión por etapa ---0,32\,GB al inicio; 0,45 tras construir la submuestra; 0,67 tras generar
las particiones; 0,88 tras el reporte de \emph{profiling}; 0,94 al cerrar el análisis no
supervisado--- es gradual y sin saltos abruptos, lo que confirma que ninguna etapa materializa el
archivo completo. La submuestra se persiste en formato \texttt{.parquet}, que ocupa 62,6\,MB en
disco frente a 113\,MB en memoria y conserva los tipos de dato al releer, garantía que el formato
CSV no ofrece.

\subsection{Particiones y ausencia de \emph{data leakage}}

Se emplea una partición estratificada 70/15/15 mediante \texttt{train\_test\_split} con
\texttt{stratify=y}, y el \texttt{StandardScaler} se ajusta \textbf{exclusivamente con el conjunto
de entrenamiento}. Validación y prueba se transforman con esos parámetros, sin re-ajuste.

La evidencia de que no hubo fuga de información reside en el contraste entre particiones
(Tabla~\ref{tab:leakage}). Train presenta media nula y desviación unitaria \emph{exactas} por
construcción, puesto que esa es la definición de la transformación; validación y prueba quedan
próximas pero no exactas, al haber sido transformadas con parámetros ajenos. Si el escalador se
hubiera ajustado sobre el conjunto completo antes de particionar, las tres filas serían exactas
por igual. La asimetría observada es precisamente la firma de un procedimiento correcto, y un
resultado exacto en las tres habría constituido el indicio de fuga.

\begin{table}[H]
\centering
\caption{Verificación del escalado sobre las tres particiones. La proporción de señal es de
52,954\,\% en las tres, efecto directo de la estratificación.}
\label{tab:leakage}
\begin{tabular}{lrrrr}
\toprule
\textbf{Partición} & \textbf{Filas} & $|\text{media}|$ \textbf{máx.} & \textbf{std mín.} & \textbf{std máx.} \\
\midrule
Train      & 699.999 & $4{,}21\times10^{-7}$ & 1,000000 & 1,000000 \\
Validación & 150.000 & $7{,}45\times10^{-3}$ & 0,992832 & 1,015338 \\
Prueba     & 150.000 & $7{,}16\times10^{-3}$ & 0,989940 & 1,005240 \\
\bottomrule
\end{tabular}
\end{table}

\subsection{Auditoría comparativa de particiones}

La auditoría exigida por el enunciado se realizó con \texttt{ydata-profiling}, empleando
\texttt{ProfileReport} con el parámetro \texttt{minimal=True} ---que desactiva el cálculo de
correlaciones costosas e interacciones--- sobre una submuestra de 10.000 filas por partición, y la
función \texttt{compare([train, val, test])}. El reporte resultante, de 3,1\,MB generado en 11
segundos, superpone las distribuciones de las 28 variables en las tres particiones: las curvas se
solapan en todos los casos, confirmando la ausencia de \emph{data drift}.

\subsection{Análisis no supervisado}

Sobre una submuestra del 10\,\% (100.000 filas) se aplicaron PCA, t-SNE \citep{tsne} y K-Means.
El análisis responde una pregunta previa al entrenamiento: ¿conforman las dos clases grupos
separables en el espacio de las variables? Si un método que no accede a las etiquetas recuperara
la separación entre señal y fondo, el problema admitiría soluciones considerablemente más simples
que una red profunda.

\begin{figure}[H]
\centering
\begin{subfigure}{0.48\textwidth}
  \includegraphics[width=\linewidth]{\figdir/pca.png}
  \caption{PCA}
\end{subfigure}\hfill
\begin{subfigure}{0.48\textwidth}
  \includegraphics[width=\linewidth]{\figdir/kmeans.png}
  \caption{K-Means}
\end{subfigure}
\caption{Análisis no supervisado sobre 100.000 observaciones. Se requieren 20 de 28 componentes
para explicar el 90\,\% de la varianza, y el agrupamiento no reproduce la etiqueta física.}
\label{fig:eda}
\end{figure}

Los tres métodos convergen en un mismo diagnóstico. El PCA requiere \textbf{20 de 28 componentes}
para explicar el 90\,\% de la varianza, y las dos primeras capturan apenas el 21,1\,\%: no existe
un subespacio lineal dominante. El índice de silueta \citep{rousseeuw} es bajo para todo $k$
evaluado ---0,186 con $k=2$ y por debajo de 0,045 para $k=3,4,5$---, y el índice de Rand ajustado
\citep{hubert} frente a la etiqueta real es de \textbf{0,0085}, prácticamente nulo. El
agrupamiento por distancia euclidiana no guarda relación alguna con la división física entre señal
y fondo: la estructura geométrica dominante está determinada por la cinemática general de las
colisiones y no por el proceso que las originó.

Lejos de constituir un resultado negativo, este diagnóstico es la justificación cuantitativa del
diseño posterior. El problema no es separable por métodos lineales ni por agrupamiento, de modo
que se requiere un modelo capaz de representar interacciones no lineales. El análisis documentó
además correlaciones de hasta 0,896 entre masas invariantes, colinealidad de origen estructural
que obliga a regularizar de forma explícita.

\section{Fase B: Ingeniería de Variables y Búsqueda de Hiperparámetros}

\subsection{Construcción de 23 variables derivadas}

Las 7 variables de alto nivel del conjunto original son, en rigor, ingeniería de variables
realizada por expertos del dominio, lo que acredita que sobre estos datos el procedimiento
funciona. El trabajo prolonga esa línea en una dirección que el conjunto no cubre: las
\textbf{relaciones angulares} entre partículas. Las masas invariantes combinan energías, pero
ninguna variable del archivo expresa el grado de separación entre dos \emph{jets}, magnitud con
significado físico, dado que la topología del evento ---si los productos emergen en direcciones
próximas u opuestas--- depende de la masa de la partícula que los originó.

La medida estándar de separación entre objetos en física de altas energías, definida en el plano
de rapidez y ángulo azimutal \citep{antikt}, es

\begin{equation}
\Delta R_{ij} = \sqrt{(\eta_i - \eta_j)^2 + (\Delta\phi_{ij})^2}
\label{eq:deltar}
\end{equation}

Se construyeron \textbf{23 variables} agrupadas en cinco familias: separación angular (seis
\texttt{dR\_*} más \texttt{deta\_j1j2} y \texttt{dphi\_j1j2}), relación leptón--resto del evento,
balance energético (\texttt{ht\_jets}, \texttt{pt\_total} y cocientes), etiquetado de sabor pesado
\citep{cms2018} mediante \texttt{n\_btags} y \texttt{suma\_btags}, y transformaciones de masa
---logaritmos, que comprimen las colas largas, y cocientes, que son adimensionales---. La
construcción completa sobre 1.000.000 de filas tomó 0,6 segundos y ninguna variable presenta
valores \texttt{NaN} ni infinitos.

\subsection{Tratamiento de la periodicidad de \texorpdfstring{$\phi$}{phi}}

El ángulo azimutal es periódico: dos objetos situados en extremos opuestos del rango son en
realidad adyacentes, por lo que una resta directa asignaría separación máxima a una configuración
de separación mínima. La diferencia debe envolverse al intervalo correspondiente a medio período:

\begin{equation}
\Delta\phi_{ij} = (\phi_i - \phi_j) - P \cdot \text{round}\!\left(\frac{\phi_i - \phi_j}{P}\right)
\label{eq:dphi}
\end{equation}

A ello se añade una particularidad del conjunto: $\phi$ \textbf{no está expresado en radianes},
sino reescalado. El período se determinó empíricamente en $P = 3{,}4879$ (rango observado
$\pm1{,}7439$), frente a los $2\pi = 6{,}2832$ que correspondería a radianes. Asumir radianes
habría falseado las siete variables angulares construidas.

\subsection{Variables más informativas}

De las 15 variables más correlacionadas con la etiqueta, \textbf{9 son construidas en esta fase}.
El hallazgo más informativo es que \texttt{m\_wbb\_over\_m\_wwbb} (0,1391) supera a
\texttt{m\_wwbb} (0,1243), que es una de las dos masas con las que se construye: la información
reside en la \emph{proporción} entre sistemas físicos y no en su magnitud absoluta, que depende de
la energía total del evento y por ello introduce variabilidad ajena a la clase. El mismo efecto se
observa en \texttt{m\_bb\_over\_m\_wbb} (0,1198) frente a \texttt{m\_wbb} (0,0663).

Cabe advertir que estos coeficientes miden correlación \emph{lineal}, mientras el análisis no
supervisado estableció que la información discriminante reside en las interacciones no lineales.
El coeficiente sirve para verificar que la construcción produjo variables con señal, no para
descartar las demás.

\subsection{Búsqueda progresiva con Optuna}

Conforme a la exigencia de no realizar la búsqueda con el 100\,\% de los datos, esta se ejecutó
sobre el \textbf{20\,\%} ---139.999 filas de entrenamiento y 30.000 de validación--- con 30
intentos en \textbf{16,9 minutos}. El muestreador TPE de Optuna \citep{optuna} modela qué regiones
del espacio produjeron buenos resultados y concentra allí los intentos siguientes, mientras el
\texttt{MedianPruner} interrumpe un intento a mitad de entrenamiento si su ROC-AUC se sitúa por
debajo de la mediana de los previos en la misma época. De los 30 intentos, \textbf{23 se
completaron y 7 fueron podados}.

Buscar sobre el 100\,\% de los datos habría demandado algo más de 80 minutos para los mismos 30
intentos: el costo medido por época sobre 700.000 filas va de 1,52\,s con la arquitectura menor a
10,16\,s con la mayor, un factor cercano a 5 respecto del subconjunto. El \emph{pruning} no fue un
refinamiento sino una condición de viabilidad en un entorno sin GPU.

El análisis de importancia atribuyó el \textbf{88,5\,\%} de la variación del objetivo al
\emph{dropout}, frente al \textbf{0,9\,\%} de la arquitectura. Leído aisladamente, ese resultado
sugeriría que la arquitectura es irrelevante; sin embargo, los diez mejores intentos del estudio
emplean \emph{todos} la arquitectura \texttt{512-256-128}, sin excepción. Ambos hechos son
compatibles: la importancia mide la varianza explicada a lo largo de todo el espacio ---y un
\emph{dropout} mal elegido degrada cualquier arquitectura, produciendo oscilaciones amplias---,
mientras el ranking muestra que alcanzar la zona superior exige la arquitectura mayor. La
regularización decide cuán bien se aprovecha una arquitectura; la arquitectura impone el techo.

\subsection{Arquitectura y motor de entrenamiento}

La red es una FFNN cuyos bloques ocultos se componen de \texttt{Linear} $\rightarrow$
\texttt{BatchNorm1d} \citep{batchnorm} $\rightarrow$ \texttt{ReLU} $\rightarrow$ \texttt{Dropout}
\citep{dropout}, con una única neurona lineal de salida que entrega \emph{logits}. Las capas
ocultas emplean inicialización He \citep{he2015}, que compensa con $\sigma = \sqrt{2/\text{fan\_in}}$
la reducción de varianza que introduce ReLU, y la capa de salida inicialización Xavier
\citep{glorot}, dado que no está seguida de una no linealidad. Las capas lineales que preceden a
\texttt{BatchNorm1d} se definen con \texttt{bias=False}, puesto que la normalización resta la
media del lote e incorpora su propio parámetro de desplazamiento.

El optimizador es AdamW \citep{adamw}, que desacopla la regularización $L_2$ del paso adaptativo
---pertinente dada la colinealidad documentada---, y la función de pérdida es
\texttt{BCEWithLogitsLoss} con

\begin{equation}
\texttt{pos\_weight} = \frac{n_{\text{neg}}}{n_{\text{pos}}} = 0{,}8884
\label{eq:posweight}
\end{equation}

El entrenamiento incorpora Early Stopping con restauración de los mejores pesos ---no los de la
última época--- y persistencia del \emph{checkpoint} en Google Drive cada vez que la métrica de
validación mejora, de modo que una interrupción no obligue a reiniciar.

\subsection{Modelos base con muestreo Bootstrap}

Cada modelo se entrena sobre una muestra Bootstrap \citep{efron} distinta, extraída con reemplazo,
que contiene alrededor del 63,2\,\% de las filas originales distintas y deja fuera el 36,8\,\%
restante (\emph{Out-Of-Bag}), conforme al esquema de Bagging \citep{breiman}. El efecto buscado es
que cada modelo observe una distribución ligeramente distinta y, en consecuencia, cometa errores
distintos.

A diferencia del esquema base, que replica una única configuración en los $M$ modelos, aquí
\textbf{cada modelo recibe una de las tres mejores configuraciones con arquitecturas distintas}
del estudio de Optuna. La razón es que modelos de configuración idéntica aprenden funciones muy
similares y sus errores quedan fuertemente correlacionados, lo que dejaría a la meta-red de la
Fase~C sin nada sustantivo que ponderar. La diversidad proviene así de tres fuentes ---la muestra
Bootstrap, la semilla de inicialización y la configuración--- en lugar de una sola.

\begin{table}[H]
\centering
\caption{Modelos base. El orden de mérito se conserva al pasar del 20\,\% al 100\,\% de los datos,
lo que valida empíricamente el esquema de búsqueda progresiva.}
\label{tab:base}
\small
\begin{tabular}{llrrrrr}
\toprule
\textbf{Modelo} & \textbf{Arquitectura} & \textbf{Dropout} & \textbf{AUC (20\,\%)} & \textbf{AUC (100\,\%)} & \textbf{Épocas} & \textbf{Min.} \\
\midrule
base\_1 & 512-256-128 & 0,302 & 0,83823 & \textbf{0,85341} & 41 & 9,0 \\
base\_2 & 256-128-64  & 0,274 & 0,83372 & 0,84993 & 45 & 3,8 \\
base\_3 & 128-64      & 0,440 & 0,82243 & 0,83332 & 43 & 1,8 \\
\bottomrule
\end{tabular}
\end{table}

Las tres muestras Bootstrap arrojaron una cobertura de 63,2\,\% de filas únicas, coincidente con
el valor teórico. Los tres modelos activaron Early Stopping entre las épocas 41 y 45, sin agotar
el presupuesto de 60, y la pérdida de validación descendió de forma monótona hasta estabilizarse,
sin el repunte que caracteriza al sobreajuste (Figura~\ref{fig:curvas}). La brecha entre pérdida
de entrenamiento y validación en la mejor época ---0,0345; 0,0214 y 0,0107 respectivamente--- es
moderada y creciente con la capacidad de la red, comportamiento esperado y contenido por la
regularización.

\begin{figure}[H]
\centering
\includegraphics[width=0.95\textwidth]{\figdir/curvas_modelos_base.png}
\caption{Curvas de pérdida (entrenamiento \emph{vs.} validación) y ROC-AUC de validación por
época. La línea vertical marca la época seleccionada por Early Stopping.}
\label{fig:curvas}
\end{figure}

\section{Fase C: Stacking Ensemble}

\subsection{Arquitectura de la meta-red}

El procedimiento más simple para unir varios modelos consiste en promediar sus predicciones con
pesos iguales, lo que supone que todos son igualmente confiables. El \emph{Stacking}
\citep{wolpert} entrena en cambio un segundo modelo ---el meta-aprendiz--- cuya única función es
determinar cómo ponderarlos, infiriendo de los datos que un modelo merece más peso que otro.

La implementación congela los $M$ modelos base, concatena sus \emph{logits} y los combina mediante
una capa lineal entrenable:

\begin{equation}
\hat{y} = \mathbf{w}^{\top} \big[ f_1(\mathbf{x}),\, f_2(\mathbf{x}),\, \ldots,\, f_M(\mathbf{x}) \big] + b
\label{eq:stacking}
\end{equation}

Los modelos base permanecen siempre en modo evaluación, incluso mientras el ensamble se entrena.
De lo contrario su \emph{dropout} se activaría y las estadísticas de \texttt{BatchNorm1d}
continuarían actualizándose, alterando precisamente las predicciones que la meta-red intenta
aprender a combinar.

\subsection{Conjunto de ajuste de la meta-red}

Si el meta-aprendiz observara las predicciones de los modelos base sobre los datos con que esos
modelos fueron entrenados, registraría un desempeño artificialmente elevado ---cada modelo ha
memorizado parcialmente esas observaciones--- y aprendería a asignarles una confianza superior a
la que les corresponde. Por ello la meta-red se entrena sobre el conjunto de \textbf{validación},
que ningún modelo base utilizó para ajustar sus pesos, dividido 70/30 para que su propio criterio
de Early Stopping no se evalúe sobre las mismas observaciones que ajusta. El conjunto de prueba
permanece sin utilizar.

\subsection{Inicialización de la meta-capa}

Se introdujo una desviación deliberada respecto del código base, que inicializa la meta-capa con
Xavier. Con solo $M=3$ entradas, esa inicialización puede entregar pesos negativos, en cuyo caso
el ensamble parte invirtiendo el ordenamiento de las predicciones ---ROC-AUC por debajo de 0,5---
y debe recorrer el espacio de parámetros hasta alcanzar el signo correcto antes de poder mejorar.
El comportamiento se verificó experimentalmente durante las pruebas. La meta-capa se inicializa
por tanto en el promedio simple, con peso $1/M$ y sesgo nulo, de modo que el ensamble parte
exactamente en la línea base que debe superar.

\section{Análisis de Resultados}

\subsection{Desempeño sobre el conjunto de prueba}

El conjunto de prueba se abrió una sola vez, al final del trabajo. No intervino en el
entrenamiento de los modelos base, ni en su Early Stopping, ni en la búsqueda de hiperparámetros,
ni en el ajuste de la meta-red, de modo que las cifras de la Tabla~\ref{tab:test} constituyen una
estimación insesgada del desempeño.

\begin{table}[H]
\centering
\caption{Desempeño sobre los 150.000 eventos del conjunto de prueba.}
\label{tab:test}
\begin{tabular}{lrrrr}
\toprule
\textbf{Predictor} & \textbf{ROC-AUC} & \textbf{PR-AUC} & \textbf{F1} & \textbf{Exactitud} \\
\midrule
Modelo base 1 (512-256-128) & 0,85316 & 0,86499 & 0,77918 & 76,92\,\% \\
Modelo base 2 (256-128-64)  & 0,84880 & 0,86094 & 0,77860 & 76,43\,\% \\
Modelo base 3 (128-64)      & 0,83279 & 0,84582 & 0,75891 & 74,97\,\% \\
Promedio simple             & 0,85282 & 0,86469 & 0,78046 & 76,89\,\% \\
\textbf{Stacking}           & \textbf{0,85572} & \textbf{0,86763} & \textbf{0,78115} & \textbf{77,13\,\%} \\
\bottomrule
\end{tabular}
\end{table}

El Stacking obtiene el mejor valor en las cuatro métricas, superando al promedio simple en 0,00290
de ROC-AUC y al mejor modelo individual en 0,00256.

El dato más relevante de la tabla, sin embargo, no es esa ganancia sino el fracaso de la
alternativa ingenua: \textbf{el promedio simple quedó por debajo del mejor modelo individual}
---0,85282 frente a 0,85316---. Combinar los tres modelos otorgándoles el mismo voto produjo un
predictor peor que quedarse únicamente con el mejor de ellos. La causa es identificable:
\texttt{base\_3} se sitúa dos centésimas de ROC-AUC por debajo de \texttt{base\_1} y un promedio
sin pesos le concede un tercio de la decisión, de modo que lo que ese tercio resta supera lo que
la diversidad de los otros dos aporta.

La ganancia absoluta del ensamble es modesta ---dos milésimas y media de ROC-AUC---, pero su
contribución no se mide contra el mejor modelo sino contra la alternativa que un ensamble ingenuo
habría producido. En un conjunto cuyos miembros no son equivalentes, aprender los pesos no
constituye un refinamiento opcional: es lo que separa mejorar de empeorar.

\subsection{Los pesos aprendidos}

\begin{table}[H]
\centering
\caption{Pesos de la meta-capa contrastados con el desempeño individual en validación.}
\label{tab:pesos}
\begin{tabular}{lrrr}
\toprule
\textbf{Modelo} & \textbf{Peso aprendido} & \textbf{Peso relativo} & \textbf{ROC-AUC individual} \\
\midrule
base\_1 & $+0{,}6401$ & 55,9\,\% & 0,85341 \\
base\_2 & $+0{,}4547$ & 39,7\,\% & 0,84993 \\
base\_3 & $-0{,}0493$ & 4,3\,\%  & 0,83332 \\
sesgo   & $-0{,}0469$ & ---      & --- \\
\bottomrule
\end{tabular}
\end{table}

La meta-red ordenó a los tres modelos \textbf{exactamente según su competencia individual} sin que
esa información se le proporcionara: solo observó sus predicciones sobre el conjunto de
validación. El modelo más preciso recibió más del doble del peso que el intermedio.

El signo negativo de \texttt{base\_3} merece atención, porque no constituye un defecto. Un peso
negativo indica que la meta-red no lo emplea como un voto adicional sino como \textbf{corrector}:
resta parte de su predicción a la combinación de los otros dos. Es el comportamiento esperable de
un modelo que subajusta ---su \emph{dropout} de 0,440 sobre la red más pequeña se lo impidió--- y
que, por lo mismo, sistematiza sus errores de un modo que resulta informativo restar. Un promedio
simple es incapaz de expresar esa relación, dado que sus pesos son positivos e iguales por
definición.

Conviene registrar lo que la meta-red no necesitó: 4 parámetros entrenables frente a 262.019
congelados de los modelos base, esto es, el 0,0015\,\% del total, ajustados en 11 segundos y 22
épocas. El Stacking no vuelve a aprender el problema, sino únicamente cómo combinar lo ya
aprendido.

\begin{figure}[H]
\centering
\begin{subfigure}{0.48\textwidth}
  \includegraphics[width=\linewidth]{\figdir/meta_stacking.png}
  \caption{Pesos aprendidos por la meta-red.}
\end{subfigure}\hfill
\begin{subfigure}{0.48\textwidth}
  \includegraphics[width=\linewidth]{\figdir/matriz_confusion.png}
  \caption{Matriz de confusión del Stacking.}
\end{subfigure}
\caption{Combinación aprendida y desempeño final sobre el conjunto de prueba.}
\label{fig:final}
\end{figure}

\subsection{Lectura en el contexto del experimento}

Con el umbral por defecto de 0,5, de las 79.431 colisiones de señal presentes en el conjunto de
prueba el modelo detectó \textbf{61.235 (77,09\,\%)} y perdió 18.196; de las 70.569 de fondo,
descartó correctamente 54.454 (77,16\,\%) y dejó pasar 16.115, que contaminan la muestra
seleccionada. De cada 100 eventos marcados como señal, \textbf{79,2 lo son efectivamente}.

Ambos errores tienen costos distintos y no intercambiables. Un falso negativo es una colisión de
Higgs descartada: tiempo de acelerador invertido en un evento que luego se ignora. Un falso
positivo contamina la muestra y diluye la significancia estadística del análisis posterior. La
simetría casi perfecta del resultado no es una propiedad deseable en sí misma, sino la
consecuencia de haber fijado el umbral en 0,5 sin un criterio físico.

El análisis de umbrales muestra el margen disponible: desplazar el corte a 0,3565 eleva el
\emph{recall} al 87,2\,\% a cambio de 6,3 puntos de precisión, mientras el valor que iguala
precisión y exhaustividad se sitúa en 0,4826. En un análisis real esa elección no se resuelve
maximizando F1, métrica agnóstica al dominio, sino en función de la significancia estadística
buscada y del presupuesto de cómputo disponible para procesar los eventos seleccionados.

\subsection{Desafíos de hardware, tiempo y memoria}

El trabajo se ejecutó \textbf{íntegramente en CPU} (PyTorch 2.15.0+cpu, 8 hilos), sin GPU
disponible, y esa restricción condicionó todas las decisiones de dimensionamiento. El tiempo total
acumulado de entrenamiento fue de aproximadamente \textbf{32 minutos} ---16,9 de búsqueda, 14,6 de
los tres modelos base y 11 segundos de la meta-red---, holgadamente dentro del presupuesto de 12 a
24 horas que fija el enunciado, holgura atribuible en buena medida al esquema de búsqueda
progresiva.

Los picos de memoria por fase fueron de 0,94\,GB en la Fase~A, 1,35\,GB en la Fase~B ---alcanzado
tras la búsqueda, que mantiene en memoria los tensores de entrenamiento y los del subconjunto--- y
0,88\,GB en la Fase~C, todos muy por debajo de los $\sim$12\,GB del nivel gratuito de Google
Colab.

Cabe señalar el bajo rendimiento marginal de la capacidad: \texttt{base\_1} necesitó 9,0 minutos
para superar a \texttt{base\_2} en 0,0035 de ROC-AUC, con 3,5 veces más parámetros y 2,4 veces más
tiempo de entrenamiento.

\section{Conclusiones y Trabajo Futuro}

El pipeline desarrollado procesa 1.000.000 de colisiones, construye 23 variables derivadas,
selecciona hiperparámetros mediante búsqueda progresiva sobre el 20\,\% de los datos, entrena tres
redes profundas sobre muestras Bootstrap distintas y las combina mediante una meta-red de pesos
entrenables, alcanzando un \textbf{ROC-AUC de 0,85572} y una \textbf{exactitud del 77,13\,\%}
sobre un conjunto de prueba de 150.000 eventos jamás utilizado durante el desarrollo.

El diagnóstico no supervisado de la Fase~A ---ARI de 0,0085 y 20 de 28 componentes necesarias para
el 90\,\% de la varianza--- anticipaba que el problema exigía un modelo capaz de representar
interacciones no lineales, y el resultado final lo confirma. La ingeniería de variables contribuyó
de forma verificable, con 9 de las 15 variables más correlacionadas con la etiqueta construidas en
la Fase~B.

El aporte metodológico más claro del trabajo se encuentra en el contraste entre el Stacking y el
promedio simple. Diversificar las configuraciones de los modelos base produjo miembros de
competencia desigual, situación en la que el promedio sin pesos degradó el resultado respecto del
mejor modelo individual mientras la combinación aprendida lo superó. Los pesos resultantes
reprodujeron el orden de competencia de los modelos y asignaron un coeficiente negativo al miembro
más débil, empleándolo como corrector en lugar de descartarlo.

Las líneas de continuación se desprenden de limitaciones documentadas y no de consideraciones
generales:

\begin{itemize}[leftmargin=1.5em,itemsep=0.3em]
  \item \textbf{Ampliar el espacio de búsqueda.} Los diez mejores intentos emplearon la
  arquitectura mayor de las cinco disponibles; cuando el óptimo queda en el borde del espacio
  explorado, es el borde y no el problema el que impone el techo. Extenderlo hacia redes más
  profundas y anchas, con más intentos e hiperparámetros adicionales ---programación de la tasa de
  aprendizaje, funciones de activación alternativas---, es la continuación natural, si bien con
  rendimientos decrecientes dado el bajo rendimiento marginal ya observado.

  \item \textbf{Contrastar la hipótesis del trabajo original.} \citet{baldi2014} sostienen que una
  red profunda puede extraer de las 21 variables de bajo nivel la información que los físicos
  codificaron en las derivadas. Este trabajo avanzó en la dirección opuesta, añadiendo variables
  construidas; la prueba simétrica ---entrenar el mismo pipeline solo con las variables de bajo
  nivel--- queda pendiente y es directamente ejecutable con el código desarrollado.

  \item \textbf{Escalar al conjunto completo.} Se empleó el 9\,\% de los datos disponibles, y hay
  evidencia de que el margen no se agota: al pasar del 20\,\% al 100\,\% del conjunto de
  entrenamiento, las tres configuraciones mejoraron entre 0,011 y 0,015 de ROC-AUC sin señales de
  saturación. Requeriría lectura incremental desde disco, al dejar los tensores de caber en
  memoria.

  \item \textbf{Diversificar por regularización.} Siendo el \emph{dropout} el hiperparámetro
  dominante, diversificar explícitamente ese eje debería producir miembros con errores aún menos
  correlacionados que diversificar la arquitectura.

  \item \textbf{Calibrar las probabilidades} y \textbf{aprovechar el conjunto Out-Of-Bag}, hoy
  descartado, para estimar el error de cada modelo base sin consumir el de validación.
\end{itemize}

% ============================================================ REFERENCIAS
\newpage
\addcontentsline{toc}{section}{Referencias}
\begin{thebibliography}{99}
\setlength{\itemsep}{0.5em}

\bibitem[Akiba et al., 2019]{optuna}
Akiba, T., Sano, S., Yanase, T., Ohta, T., \& Koyama, M. (2019). Optuna: A Next-generation
Hyperparameter Optimization Framework. \emph{Proceedings of the 25th ACM SIGKDD International
Conference on Knowledge Discovery \& Data Mining}, 2623-2631.
\url{https://doi.org/10.1145/3292500.3330701}

\bibitem[Baldi et al., 2014]{baldi2014}
Baldi, P., Sadowski, P., \& Whiteson, D. (2014). Searching for exotic particles in high-energy
physics with deep learning. \emph{Nature Communications}, \emph{5}, 4308.
\url{https://doi.org/10.1038/ncomms5308}

\bibitem[Breiman, 1996]{breiman}
Breiman, L. (1996). Bagging Predictors. \emph{Machine Learning}, \emph{24}(2), 123-140.
\url{https://doi.org/10.1007/BF00058655}

\bibitem[Cacciari et al., 2008]{antikt}
Cacciari, M., Salam, G. P., \& Soyez, G. (2008). The anti-$k_t$ jet clustering algorithm.
\emph{Journal of High Energy Physics}, \emph{2008}(04), 063.
\url{https://doi.org/10.1088/1126-6708/2008/04/063}

\bibitem[CMS Collaboration, 2018]{cms2018}
CMS Collaboration. (2018). Identification of heavy-flavour jets with the CMS detector in pp
collisions at 13 TeV. \emph{Journal of Instrumentation}, \emph{13}, P05011.
\url{https://doi.org/10.1088/1748-0221/13/05/P05011}

\bibitem[Efron, 1979]{efron}
Efron, B. (1979). Bootstrap Methods: Another Look at the Jackknife. \emph{The Annals of
Statistics}, \emph{7}(1), 1-26. \url{https://doi.org/10.1214/aos/1176344552}

\bibitem[Glorot \& Bengio, 2010]{glorot}
Glorot, X., \& Bengio, Y. (2010). Understanding the difficulty of training deep feedforward neural
networks. \emph{Proceedings of AISTATS}, PMLR 9, 249-256.
\url{https://proceedings.mlr.press/v9/glorot10a.html}

\bibitem[He et al., 2015]{he2015}
He, K., Zhang, X., Ren, S., \& Sun, J. (2015). Delving Deep into Rectifiers: Surpassing
Human-Level Performance on ImageNet Classification. \emph{Proceedings of ICCV}, 1026-1034.
\url{https://arxiv.org/abs/1502.01852}

\bibitem[Hubert \& Arabie, 1985]{hubert}
Hubert, L., \& Arabie, P. (1985). Comparing partitions. \emph{Journal of Classification},
\emph{2}(1), 193-218. \url{https://doi.org/10.1007/BF01908075}

\bibitem[Ioffe \& Szegedy, 2015]{batchnorm}
Ioffe, S., \& Szegedy, C. (2015). Batch Normalization: Accelerating Deep Network Training by
Reducing Internal Covariate Shift. \emph{Proceedings of ICML}, PMLR 37, 448-456.
\url{https://proceedings.mlr.press/v37/ioffe15.html}

\bibitem[Loshchilov \& Hutter, 2019]{adamw}
Loshchilov, I., \& Hutter, F. (2019). Decoupled Weight Decay Regularization. \emph{International
Conference on Learning Representations (ICLR)}. \url{https://arxiv.org/abs/1711.05101}

\bibitem[Rousseeuw, 1987]{rousseeuw}
Rousseeuw, P. J. (1987). Silhouettes: a graphical aid to the interpretation and validation of
cluster analysis. \emph{Journal of Computational and Applied Mathematics}, \emph{20}, 53-65.
\url{https://doi.org/10.1016/0377-0427(87)90125-7}

\bibitem[Srivastava et al., 2014]{dropout}
Srivastava, N., Hinton, G., Krizhevsky, A., Sutskever, I., \& Salakhutdinov, R. (2014). Dropout: A
Simple Way to Prevent Neural Networks from Overfitting. \emph{Journal of Machine Learning
Research}, \emph{15}, 1929-1958. \url{https://jmlr.org/papers/v15/srivastava14a.html}

\bibitem[van der Maaten \& Hinton, 2008]{tsne}
van der Maaten, L., \& Hinton, G. (2008). Visualizing Data using t-SNE. \emph{Journal of Machine
Learning Research}, \emph{9}, 2579-2605. \url{https://jmlr.org/papers/v9/vandermaaten08a.html}

\bibitem[Whiteson, 2014]{uci}
Whiteson, D. (2014). \emph{HIGGS} [Conjunto de datos]. UCI Machine Learning Repository, ID 280.
\url{https://archive.ics.uci.edu/dataset/280/higgs}

\bibitem[Wolpert, 1992]{wolpert}
Wolpert, D. H. (1992). Stacked generalization. \emph{Neural Networks}, \emph{5}, 241-259.

\end{thebibliography}

\end{document}
